% ChargeShield-FL --- DSN 2027 --- 7. Conclusion
% Scritta il 2026-09-23. Solo risposte alle RQ effettivamente affrontate
% (guida sez. 10): per ora e' una conclusione di stato, da riscrivere dopo
% E-A, E-B, E-E, E-D.

\section{Conclusion}
\label{sec:conclusion}

We set out to establish under which conditions differential privacy reduces
the ability of membership inference attacks in a federated autoencoder trained
on real EV-charging sessions, and at what cost. The evidence verified so far
answers part of RQ1. Without DP, average attack metrics are at chance while a
per-record analysis finds a systematic excess of recognisable sessions; with
client-level DP at $\varepsilon \le 1$ the excess disappears together with the
usefulness of the model, so that regime does not tell whether DP protects. A
canary control shows that the instrument detects memorisation when it is
induced on purpose. \pending{Answer to RQ1 from the $\varepsilon$ sweep and
the record-level comparison; answers to RQ3 and RQ2, or an explicit statement
that they were deferred, with the reason and its consequence on the claim.}
